\documentclass[12pt,a4paper,oneside]{report}

% ===================================================================
% PACKAGES ESSENTIELS
% ===================================================================

% Encodage et langue
\usepackage[utf8]{inputenc}
\usepackage[T1]{fontenc}
\usepackage[french]{babel}

% Mise en page
\usepackage[top=2.5cm, bottom=2.5cm, left=3cm, right=2.5cm]{geometry}
\usepackage{setspace}  % Interligne
\onehalfspacing  % Interligne 1.5

% Math et symboles
\usepackage{amsmath, amssymb, amsthm}

% Graphiques et tableaux
\usepackage{graphicx}
\usepackage{float}
\usepackage{caption}
\usepackage{subcaption}
\usepackage{booktabs}  % Tableaux professionnels
\usepackage{multirow}
\usepackage{array}

% Couleurs et liens
\usepackage[dvipsnames]{xcolor}
\usepackage[hidelinks]{hyperref}  % Liens cliquables sans encadrement
\hypersetup{
    colorlinks=true,
    linkcolor=blue,
    citecolor=blue,
    urlcolor=blue
}

% Bibliographie
\usepackage[backend=biber, style=numeric, sorting=none]{biblatex}
\addbibresource{bibliographie.bib}

% Code source (si nécessaire)
\usepackage{listings}
\lstset{
    basicstyle=\ttfamily\small,
    breaklines=true,
    frame=single,
    numbers=left,
    numberstyle=\tiny,
    keywordstyle=\color{blue},
    commentstyle=\color{green!60!black},
    stringstyle=\color{red}
}

% Algorithmes
\usepackage{algorithm}
\usepackage{algpseudocode}

% Glossaire et acronymes
\usepackage[acronym]{glossaries}
\makeglossaries

% Chapitres et sections
\usepackage{titlesec}

% En-têtes et pieds de page
\usepackage{fancyhdr}
\pagestyle{fancy}
\fancyhf{}
\fancyhead[L]{\nouppercase{\leftmark}}
\fancyhead[R]{\thepage}
\renewcommand{\headrulewidth}{0.5pt}

% ===================================================================
% INFORMATIONS DU DOCUMENT
% ===================================================================

\title{Détection précoce du Diabète de Type 1 en Afrique subsaharienne par apprentissage automatique : identification de biomarqueurs génétiques adaptés aux populations camerounaises}

\author{Sorelle}

\date{Novembre 2025}

% ===================================================================
% DÉBUT DU DOCUMENT
% ===================================================================

\begin{document}

% Page de garde
\begin{titlepage}
    \centering

    % Logo université (à ajouter)
    % \includegraphics[width=0.3\textwidth]{figures/logo_universite.png}\\[1cm]

    {\Large \textbf{[NOM DE L'UNIVERSITÉ]}}\\[0.5cm]
    {\large [Faculté / Département]}\\[0.3cm]
    {\large Master 2 Physique Atomique, Moléculaire et Biophysique}\\[2cm]

    % Titre
    {\huge \textbf{Détection précoce du Diabète de Type 1 en Afrique subsaharienne par apprentissage automatique}}\\[0.5cm]

    {\LARGE Identification de biomarqueurs génétiques adaptés aux populations camerounaises}\\[3cm]

    % Auteur
    {\Large \textbf{Présenté par :}}\\[0.3cm]
    {\Large Sorelle}\\[2cm]

    % Encadrement
    \begin{minipage}{0.8\textwidth}
        \begin{flushleft}
            {\large \textbf{Encadrant :}}\\
            {\large Dr. TCHAPET NJAFA Jean-Pierre}\\[0.5cm]

            {\large \textbf{Conseiller scientifique :}}\\
            {\large Prof. NANA Engo}
        \end{flushleft}
    \end{minipage}\\[2cm]

    \vfill

    % Date et contexte
    {\large Mémoire de Master 2}\\[0.2cm]
    {\large En vue de l'obtention du diplôme de Master}\\[0.2cm]
    {\large Soutenance publique}\\[1cm]

    {\Large \textbf{Novembre 2025}}\\[0.5cm]

    {\large Yaoundé, Cameroun}

\end{titlepage}


% Page blanche
\thispagestyle{empty}
\clearpage
\null
\clearpage

% Résumés
\pagenumbering{roman}
\setcounter{page}{1}

% Résumé en français
\chapter*{Résumé}
\addcontentsline{toc}{chapter}{Résumé}

Le Diabète de Type 1 (DT1) en Afrique subsaharienne présente des spécificités cliniques et génétiques distinctes des formes observées en Occident. Les formes idiopathiques, caractérisées par l'absence d'auto-anticorps, représentent 20 à 60\% des cas, rendant le diagnostic précoce particulièrement difficile. Cette problématique est aggravée par le \textit{dataset shift} : les modèles d'apprentissage automatique entraînés sur des populations caucasiennes montrent des performances dégradées lorsqu'appliqués aux populations africaines.

\vspace{0.3cm}

Ce mémoire de Master 2 vise à développer un système de détection précoce du DT1 spécifiquement adapté aux populations camerounaises, en utilisant des techniques d'apprentissage automatique pour identifier des biomarqueurs génétiques pertinents. Trois biomarqueurs candidats (ANP32A-IT1, ESCO2, NBPF1) ont été sélectionnés sur la base de la littérature scientifique récente montrant leur potentiel diagnostique (AUC > 0.8).

\vspace{0.3cm}

Notre méthodologie repose sur un dataset synthétique réaliste de 1000 patients camerounais, comprenant des données démographiques, cliniques et génétiques. Six algorithmes de classification supervisée ont été comparés (Régression Logistique, Arbre de Décision, Random Forest, XGBoost, LightGBM, SVM), avec optimisation des hyperparamètres par recherche en grille et validation croisée stratifiée. L'interprétabilité des prédictions a été assurée par l'utilisation de SHAP (\textit{SHapley Additive exPlanations}) et LIME (\textit{Local Interpretable Model-agnostic Explanations}).

\vspace{0.3cm}

Les résultats montrent que [À COMPLÉTER AVEC VOS RÉSULTATS]. Le modèle final atteint un recall de [XX\%] et une précision de [XX\%] sur l'ensemble de test, avec une AUC-ROC de [0.XX]. L'analyse SHAP confirme l'importance des biomarqueurs ESCO2 (importance relative : [XX\%]) et ANP32A-IT1 ([XX\%]) dans les prédictions, en cohérence avec les mécanismes physiopathologiques connus du DT1.

\vspace{0.3cm}

Une application de démonstration interactive a été développée avec Streamlit, permettant aux cliniciens d'uploader des données patients, d'obtenir des prédictions en temps réel, et de visualiser l'importance des biomarqueurs pour chaque diagnostic. Cette interface facilite le transfert de connaissance entre recherche et pratique clinique.

\vspace{0.3cm}

\textbf{Mots-clés :} Diabète de Type 1, Afrique subsaharienne, Cameroun, Apprentissage automatique, Biomarqueurs génétiques, SHAP, Santé publique

\clearpage

% Abstract en anglais
\chapter*{Abstract}
\addcontentsline{toc}{chapter}{Abstract}

Type 1 Diabetes (T1D) in sub-Saharan Africa presents clinical and genetic specificities distinct from forms observed in Western populations. Idiopathic forms, characterized by the absence of autoantibodies, represent 20 to 60\% of cases, making early diagnosis particularly challenging. This issue is exacerbated by dataset shift: machine learning models trained on Caucasian populations show degraded performance when applied to African populations.

\vspace{0.3cm}

This Master's thesis aims to develop an early detection system for T1D specifically adapted to Cameroonian populations, using machine learning techniques to identify relevant genetic biomarkers. Three candidate biomarkers (ANP32A-IT1, ESCO2, NBPF1) were selected based on recent scientific literature showing their diagnostic potential (AUC > 0.8).

\vspace{0.3cm}

Our methodology relies on a realistic synthetic dataset of 1000 Cameroonian patients, including demographic, clinical, and genetic data. Six supervised classification algorithms were compared (Logistic Regression, Decision Tree, Random Forest, XGBoost, LightGBM, SVM), with hyperparameter optimization through grid search and stratified cross-validation. Prediction interpretability was ensured using SHAP (SHapley Additive exPlanations) and LIME (Local Interpretable Model-agnostic Explanations).

\vspace{0.3cm}

Results show that [TO BE COMPLETED WITH YOUR RESULTS]. The final model achieves a recall of [XX\%] and precision of [XX\%] on the test set, with an AUC-ROC of [0.XX]. SHAP analysis confirms the importance of biomarkers ESCO2 (relative importance: [XX\%]) and ANP32A-IT1 ([XX\%]) in predictions, consistent with known T1D pathophysiological mechanisms.

\vspace{0.3cm}

An interactive demonstration application was developed with Streamlit, allowing clinicians to upload patient data, obtain real-time predictions, and visualize biomarker importance for each diagnosis. This interface facilitates knowledge transfer between research and clinical practice.

\vspace{0.3cm}

\textbf{Keywords:} Type 1 Diabetes, Sub-Saharan Africa, Cameroon, Machine Learning, Genetic Biomarkers, SHAP, Public Health

\clearpage

% Remerciements
\chapter*{Remerciements}
\addcontentsline{toc}{chapter}{Remerciements}

Je tiens à exprimer ma profonde gratitude à toutes les personnes qui ont contribué, de près ou de loin, à la réalisation de ce mémoire de Master 2.

\vspace{0.3cm}

En premier lieu, je remercie chaleureusement mon encadrant, le \textbf{Dr. TCHAPET NJAFA Jean-Pierre}, pour son accompagnement rigoureux, sa disponibilité constante et ses conseils avisés tout au long de ces quatorze semaines de travail. Sa expertise en biophysique et en intelligence artificielle médicale a été déterminante dans l'orientation et la qualité de ce projet.

\vspace{0.3cm}

Mes remerciements vont également au \textbf{Professeur NANA Engo} pour ses remarques constructives et ses orientations lors de l'affinement du sujet de recherche. Ses connaissances approfondies sur le contexte de la santé publique au Cameroun ont enrichi considérablement la pertinence de cette étude.

\vspace{0.3cm}

Je remercie [À COMPLÉTER : membres du jury, enseignants du Master, laboratoire d'accueil, etc.].

\vspace{0.3cm}

Sur le plan scientifique, je souhaite saluer les travaux pionniers des Professeurs Jean-Claude Mbanya, Eugene Sobngwi et Andre Pascal Kengne, dont les publications sur le diabète en Afrique subsaharienne ont constitué les fondations théoriques de cette recherche.

\vspace{0.3cm}

Je remercie également [À COMPLÉTER : collègues de promotion, famille, proches] pour leur soutien moral et leurs encouragements durant cette période intensive de travail.

\vspace{0.3cm}

Enfin, je dédie ce travail aux millions de personnes atteintes de diabète en Afrique subsaharienne, avec l'espoir que la recherche scientifique contribue à améliorer leur accès au diagnostic précoce et aux soins de qualité.

\clearpage

% Table des matières
\tableofcontents
\clearpage

% Liste des figures
\listoffigures
\addcontentsline{toc}{chapter}{Liste des figures}
\clearpage

% Liste des tableaux
\listoftables
\addcontentsline{toc}{chapter}{Liste des tableaux}
\clearpage

% Liste des acronymes (à compléter)
\chapter*{Liste des acronymes}
\addcontentsline{toc}{chapter}{Liste des acronymes}

\begin{description}
    \item[AUC] Area Under the Curve (Aire sous la courbe)
    \item[CNN] Convolutional Neural Network (Réseau de neurones convolutif)
    \item[CSV] Comma-Separated Values (Valeurs séparées par des virgules)
    \item[DT1] Diabète de Type 1
    \item[DT2] Diabète de Type 2
    \item[EDA] Exploratory Data Analysis (Analyse exploratoire des données)
    \item[FN] False Negative (Faux Négatif)
    \item[FP] False Positive (Faux Positif)
    \item[HbA1c] Hémoglobine glyquée
    \item[HLA] Human Leukocyte Antigen (Antigène des leucocytes humains)
    \item[IA] Intelligence Artificielle
    \item[IMC] Indice de Masse Corporelle
    \item[K-NN] K-Nearest Neighbors (K plus proches voisins)
    \item[LIME] Local Interpretable Model-agnostic Explanations
    \item[ML] Machine Learning (Apprentissage automatique)
    \item[OMS] Organisation Mondiale de la Santé
    \item[PDF] Probability Density Function (Fonction de densité de probabilité)
    \item[RAM] Random Access Memory (Mémoire vive)
    \item[RF] Random Forest (Forêt aléatoire)
    \item[ROC] Receiver Operating Characteristic
    \item[SHAP] SHapley Additive exPlanations
    \item[SMOTE] Synthetic Minority Over-sampling Technique
    \item[SVM] Support Vector Machine (Machine à vecteurs de support)
    \item[TN] True Negative (Vrai Négatif)
    \item[TP] True Positive (Vrai Positif)
    \item[XGBoost] Extreme Gradient Boosting
\end{description}

\clearpage

% Corps du mémoire
\pagenumbering{arabic}
\setcounter{page}{1}

% Chapitre 1 : Introduction
\chapter{Introduction}
\label{chap:introduction}

\section{Contexte et problématique}

Le Diabète de Type 1 (DT1) constitue un enjeu majeur de santé publique à l'échelle mondiale, avec une incidence croissante estimée à 3\% par an \cite{DiabetesAtlas2021}. En Afrique subsaharienne, cette pathologie présente des spécificités cliniques, génétiques et épidémiologiques qui la distinguent significativement des formes observées dans les populations caucasiennes.

\subsection{Le diabète de type 1 en Afrique subsaharienne}

Contrairement aux idées reçues, le DT1 n'est pas rare en Afrique subsaharienne. Des études récentes menées au Cameroun, en Ouganda et au Kenya rapportent une prévalence non négligeable, bien que probablement sous-estimée en raison des difficultés diagnostiques et du manque d'accès aux soins \cite{Mbanya2010,Katte2023}. La mortalité infantile associée au DT1 demeure dramatiquement élevée dans la région, avec des taux de survie à 5 ans inférieurs à 50\% dans certaines zones rurales, principalement dus à un diagnostic tardif et à un accès limité à l'insuline \cite{Atun2017}.

Une particularité majeure du DT1 en Afrique réside dans la fréquence élevée des \textbf{formes idiopathiques}, caractérisées par une perte de sécrétion d'insuline endogène sans signes d'auto-immunité détectables (absence d'auto-anticorps anti-îlots) \cite{Sobngwi2002}. Ces formes représentent 20 à 60\% des cas selon les études, contre moins de 5\% dans les populations européennes. Cette observation suggère des mécanismes physiopathologiques distincts, possiblement liés à des facteurs génétiques et environnementaux spécifiques aux populations africaines.

\subsection{Le problème du \textit{dataset shift}}

L'essor de l'intelligence artificielle et de l'apprentissage automatique (Machine Learning, ML) a révolutionné de nombreux domaines médicaux, incluant le diagnostic précoce de pathologies complexes. Cependant, un défi majeur émerge lorsque l'on tente d'appliquer ces modèles à des populations différentes de celles sur lesquelles ils ont été entraînés : le \textbf{\textit{dataset shift}} \cite{Quinonero-Candela2009}.

Ce phénomène se manifeste particulièrement dans le contexte du DT1 : les modèles prédictifs développés sur des cohortes américaines ou européennes (majoritairement caucasiennes) montrent des performances dégradées lorsqu'appliqués aux populations africaines. Une étude récente a démontré qu'un score de risque génétique (GRS) entraîné sur des populations européennes présentait une sensibilité de seulement 53\% sur des populations camerounaises et ougandaises, contre 91\% sur des populations européennes \cite{Oram2025}. Cette différence s'explique par :

\begin{itemize}
    \item Des \textbf{distributions génétiques} différentes : les haplotypes HLA associés au DT1 varient significativement entre populations africaines et caucasiennes.
    \item Des \textbf{facteurs environnementaux} distincts : nutrition, infections, microbiome intestinal.
    \item Des \textbf{présentations cliniques} hétérogènes : formes idiopathiques vs. auto-immunes.
\end{itemize}

\subsection{Biomarqueurs génétiques prometteurs}

Des travaux récents en bioinformatique ont identifié trois biomarqueurs génétiques candidats présentant un fort potentiel diagnostique pour le DT1 \cite{Zhang2022} :

\begin{description}
    \item[ANP32A-IT1 :] Long ARN non codant impliqué dans la régulation de la réponse immunitaire, notamment l'activation des lymphocytes T. Son expression est significativement augmentée chez les patients DT1 (AUC > 0.8).

    \item[ESCO2 :] Gène codant pour une protéine impliquée dans la régulation du cycle cellulaire et la cohésion des chromatides s\oe urs. Des mutations de ce gène sont associées à des dysfonctions pancréatiques et à l'apoptose des cellules bêta.

    \item[NBPF1 :] Membre de la famille des gènes NBPF (\textit{Neuroblastoma Breakpoint Family}), exprimé dans les cellules pancréatiques et impliqué dans la prolifération cellulaire. Son niveau d'expression corrèle avec le risque de développement du DT1.
\end{description}

Ces biomarqueurs ont été validés sur des cohortes asiatiques et européennes, mais leur pertinence dans le contexte africain reste à établir. L'adaptation de leur utilisation aux spécificités génétiques des populations camerounaises constitue une nécessité scientifique et clinique.

\section{Objectifs de la recherche}

\subsection{Objectif principal}

L'objectif principal de ce mémoire est de \textbf{développer un système de détection précoce du Diabète de Type 1 adapté aux populations camerounaises}, en utilisant des techniques d'apprentissage automatique pour identifier et valider des biomarqueurs génétiques pertinents.

\subsection{Objectifs spécifiques}

Pour atteindre cet objectif principal, nous nous fixons les objectifs spécifiques suivants :

\begin{enumerate}
    \item \textbf{Constituer un dataset synthétique réaliste} de patients camerounais intégrant des données démographiques, cliniques et génétiques, en s'appuyant sur les distributions observées dans la littérature scientifique africaine.

    \item \textbf{Comparer les performances de six algorithmes de classification supervisée} (Régression Logistique, Arbre de Décision, Random Forest, XGBoost, LightGBM, SVM) pour la prédiction du DT1, en optimisant leurs hyperparamètres par validation croisée.

    \item \textbf{Identifier et valider biologiquement les biomarqueurs génétiques les plus pertinents} parmi ANP32A-IT1, ESCO2 et NBPF1, en utilisant des techniques d'interprétabilité (SHAP, LIME).

    \item \textbf{Développer une application de démonstration interactive} permettant aux cliniciens d'utiliser le modèle pour des prédictions en temps réel, avec visualisation de l'importance des biomarqueurs.

    \item \textbf{Évaluer les performances du système} en termes de recall (sensibilité), precision et AUC-ROC, avec une priorité clinique sur la minimisation des faux négatifs.
\end{enumerate}

\section{Hypothèses de recherche}

Nos travaux reposent sur les hypothèses suivantes :

\begin{itemize}
    \item \textbf{H1 :} Les biomarqueurs ANP32A-IT1, ESCO2 et NBPF1 présentent une capacité discriminante significative (AUC > 0.70) pour la détection du DT1 dans le contexte camerounais.

    \item \textbf{H2 :} Un modèle d'apprentissage automatique entraîné sur des données de populations africaines surpasse les approches basées sur des cohortes occidentales en termes de sensibilité et de spécificité.

    \item \textbf{H3 :} Les algorithmes d'ensemble (Random Forest, XGBoost, LightGBM) offrent de meilleures performances que les modèles linéaires pour capturer les interactions complexes entre biomarqueurs.

    \item \textbf{H4 :} L'interprétabilité des prédictions via SHAP permet de valider biologiquement l'importance des biomarqueurs identifiés, renforçant la confiance clinique dans le système.
\end{itemize}

\section{Contributions scientifiques}

Ce travail apporte plusieurs contributions originales :

\begin{itemize}
    \item \textbf{Contribution méthodologique :} Développement d'un pipeline reproductible de détection précoce du DT1 adapté aux contraintes de ressources computationnelles limitées (pas de GPU, RAM modérée), typiques du contexte camerounais.

    \item \textbf{Contribution scientifique :} Première étude combinant les biomarqueurs ANP32A-IT1, ESCO2 et NBPF1 dans le contexte spécifique des populations camerounaises, avec validation de leur pertinence par des méthodes d'interprétabilité avancées.

    \item \textbf{Contribution pratique :} Création d'un outil de démonstration accessible (application web Streamlit) facilitant le transfert de connaissances entre recherche et pratique clinique, utilisable sans expertise technique approfondie.

    \item \textbf{Contribution pédagogique :} Documentation exhaustive du projet (notebooks Jupyter commentés, guides d'installation, fiches conceptuelles) favorisant la reproductibilité et l'appropriation des méthodes par d'autres chercheurs africains.
\end{itemize}

\section{Organisation du mémoire}

Ce mémoire s'organise en six chapitres :

\begin{description}
    \item[Chapitre 1 - Introduction :] Contextualisation de la problématique du DT1 en Afrique subsaharienne, présentation du \textit{dataset shift}, justification du choix des biomarqueurs et énoncé des objectifs.

    \item[Chapitre 2 - État de l'art :] Revue de la littérature sur trois axes : (i) biologie du DT1 et spécificités africaines, (ii) intelligence artificielle en médecine et métriques d'évaluation, (iii) applications existantes du ML pour le DT1 et leurs limitations.

    \item[Chapitre 3 - Méthodologie :] Description détaillée du dataset synthétique, des algorithmes de classification testés, des techniques d'optimisation des hyperparamètres, des métriques d'évaluation privilégiées, et des outils d'interprétabilité utilisés.

    \item[Chapitre 4 - Résultats :] Présentation des performances des modèles, comparaisons statistiques, analyse de l'importance des biomarqueurs via SHAP, et démonstration de l'application Streamlit.

    \item[Chapitre 5 - Discussion :] Interprétation des résultats au regard de la littérature, validation biologique des biomarqueurs identifiés, analyse des limitations de l'étude, et implications pour la santé publique camerounaise.

    \item[Chapitre 6 - Conclusion :] Synthèse des contributions, réponse aux hypothèses de recherche, perspectives d'amélioration et travaux futurs.
\end{description}

\section{Délimitation du champ de l'étude}

Il est important de préciser les limites de cette recherche :

\begin{itemize}
    \item \textbf{Population cible :} Ce travail se concentre sur les populations camerounaises. La généralisabilité à d'autres pays d'Afrique subsaharienne nécessite une validation externe.

    \item \textbf{Nature des données :} Le dataset utilisé est synthétique, généré à partir de distributions réalistes extraites de la littérature. Une validation sur données réelles camerounaises constitue une étape ultérieure indispensable.

    \item \textbf{Biomarqueurs :} Nous nous focalisons sur trois biomarqueurs candidats (ANP32A-IT1, ESCO2, NBPF1). D'autres marqueurs génétiques, épigénétiques ou cliniques pourraient enrichir le modèle.

    \item \textbf{Portée clinique :} Ce système constitue un \textbf{outil d'aide à la décision}, non un remplacement du diagnostic médical. Toute prédiction positive doit être confirmée par des examens cliniques et biologiques complémentaires.

    \item \textbf{Validation réglementaire :} L'application développée est une démonstration de recherche. Son déploiement en conditions réelles nécessiterait des validations cliniques, éthiques et réglementaires approfondies.
\end{itemize}

Ces délimitations assurent une interprétation rigoureuse des résultats et guident les travaux futurs nécessaires pour un impact clinique effectif.


% Chapitre 2 : État de l'art
\chapter{État de l'art}
\label{chap:etat_art}

\section{Biologie du Diabète de Type 1}

\subsection{Physiopathologie générale}
Le DT1 est caractérisé par une destruction des cellules bêta pancréatiques productrices d'insuline. [À COMPLÉTER : mécanismes auto-immuns, génétique HLA, facteurs environnementaux]

\subsection{Spécificités africaines}
Formes idiopathiques (20-60\%), absence d'auto-anticorps, facteurs génétiques distincts \cite{Katte2023,Sobngwi2002}.

\section{Intelligence Artificielle en médecine}

\subsection{Machine Learning supervisé}
Classification, régression, métriques (recall, precision, AUC-ROC) \cite{Bishop2006,Hastie2009}.

\subsection{Algorithmes de classification}
Random Forest \cite{Breiman2001}, XGBoost \cite{Chen2016}, LightGBM \cite{Ke2017}, SVM.

\subsection{Interprétabilité}
SHAP \cite{Lundberg2017}, LIME \cite{Ribeiro2016}. Importance pour confiance clinique.

\section{Applications ML pour le DT1}

\subsection{Modèles existants}
Revue des systèmes de prédiction, scores de risque génétique \cite{Contreras2018}.

\subsection{Limitations en Afrique}
Dataset shift \cite{Quinonero-Candela2009}, manque de données locales, performances dégradées \cite{Oram2025}.

\section{Positionnement de notre approche}
Notre travail se distingue par : (1) focus populations camerounaises, (2) biomarqueurs validés, (3) interprétabilité prioritaire.


% Chapitre 3 : Méthodologie
\chapter{Méthodologie}
\label{chap:methodologie}

% ============================================================
\section{Acquisition et préparation des données}
\label{sec:donnees}
% ============================================================

La phase de préparation des données constitue l'étape la plus critique du pipeline d'apprentissage
automatique, absorbant typiquement 80\,\% de l'effort total du projet \cite{sorelle2025demarche}.
Elle s'articule autour du processus KDD (\textit{Knowledge Discovery in Databases}) et du modèle
industriel CRISP-DM, garantissant la transition de données cliniques brutes vers un dataset raffiné
et prêt pour la modélisation \cite{contreras2018ai, dagliati2018ml}.

% -------------------------------------------------------
\subsection{Description du dataset : source, taille et variables}
\label{subsec:dataset}
% -------------------------------------------------------

Les données proviennent de l'étude multicentrique internationale \textbf{YODA}
(\textit{Young-Onset Diabetes in sub-Saharan Africa}) \cite{katte2025lancet, oram2025ancestry}.

\paragraph{Source et population.}
Le recrutement a été effectué dans des centres cliniques au Cameroun (Yaoundé et Bafoussam),
en Ouganda et en Afrique du Sud. La cohorte cible des individus d'ascendance africaine
diagnostiqués cliniquement avant l'âge de 30 ans \cite{katte2025lancet}.

\paragraph{Taille de l'échantillon.}
La population initiale est de \textbf{1\,072 participants}. Après les étapes de filtrage détaillées
ci-dessous, la cohorte finale retenue est de \textbf{894 individus} (454 hommes et 440 femmes)
— soit un taux de rétention de 83,4\,\% \cite{katte2025lancet}.

\paragraph{Variables disponibles.}
Le dataset inclut quatre catégories de variables :
\begin{itemize}
    \item \textbf{Paramètres anthropométriques} : âge au diagnostic, indice de masse corporelle (IMC\,kg/m\textsuperscript{2}).
    \item \textbf{Paramètres métaboliques} : peptide~C (nmol/L), hémoglobine glyquée HbA1c\,(\%), glycémie à jeun (mg/dL).
    \item \textbf{Marqueurs immunologiques} : auto-anticorps GADA, IA-2A, ZnT8A (positif/négatif).
    \item \textbf{Biomarqueurs génomiques (hub genes)} : expressions relatives des gènes
          ANP32A-IT1, ESCO2 et NBPF1, identifiés par analyse de réseaux de co-expression
          comme prédicteurs potentiels du DT1 \cite{wang2023genemodel}.
\end{itemize}

La variable cible est le \textbf{statut DT1} (positif = 1, négatif = 0), avec un taux de
prévalence initial de 25\,\% dans le dataset synthétique utilisé pour la modélisation.

% -------------------------------------------------------
\subsection{Méthodes de nettoyage appliquées}
\label{subsec:nettoyage}
% -------------------------------------------------------

Le nettoyage vise à réduire le bruit clinique et à assurer la stabilité des futurs modèles.
Les décisions ont été prises selon les seuils et protocoles suivants \cite{dagliati2018ml,
sorelle2025demarche} :

\paragraph{Filtrage phénotypique.}
Quatre-vingt-dix individus présentant un IMC $\geq 30\,\text{kg/m}^2$ ont été exclus afin
d'écarter les risques de confusion diagnostique avec le diabète de type 2 (DT2). De plus,
88 participants ont été retirés en l'absence d'échantillons de sérum pour la mesure
des auto-anticorps \cite{katte2025lancet}.

\paragraph{Gestion des données manquantes.}
Les taux de données manquantes observés dans l'état brut de la cohorte YODA sont présentés
dans le tableau~\ref{tab:missing}. Les variables dépassant le seuil de 30\,\% de valeurs
manquantes (triglycérides, cholestérol total) ont été supprimées du modèle final.
Pour les variables restantes, l'imputation a été réalisée via l'algorithme
\textbf{missForest} (forêt aléatoire itérative), réduisant l'erreur quadratique
moyenne (RMSE) sur l'IMC de 3{,}23 à 0{,}57 — performance significativement
supérieure à une simple imputation par la moyenne \cite{dagliati2018ml}.

\begin{table}[h!]
\centering
\caption{Taux de données manquantes observés avant nettoyage (cohorte YODA brute)}
\label{tab:missing}
\begin{tabular}{lcc}
\hline
\textbf{Variable} & \textbf{Taux manquant (\%)} & \textbf{Action appliquée} \\
\hline
HbA1c              & 16{,}9                   & Imputation missForest \\
Cholestérol total  & 34{,}1                   & Suppression de la variable \\
Triglycérides      & 36{,}2                   & Suppression de la variable \\
\hline
\end{tabular}
\end{table}

\paragraph{Traitement des valeurs aberrantes.}
Les outliers ont été traités par \textbf{Winsorization} (écrêtage) aux seuils
$[Q_1 - 1{,}5 \times IQR \;;\; Q_3 + 1{,}5 \times IQR]$ pour les variables :
ANP32A\_IT1, ESCO2, NBPF1, glycémie à jeun, HbA1c, IMC et âge. Cette méthode préserve
toute la population (aucune ligne supprimée) tout en atténuant l'influence des
valeurs extrêmes sur les paramètres des algorithmes.

\paragraph{Normalisation génomique.}
Les expressions d'ARNm des hub genes ont été traitées par l'algorithme
\textbf{RMA} (\textit{Robust Multi-array Average}) et soumises à une correction
des effets de lot via \textbf{ComBat}, garantissant la comparabilité des mesures
inter-centres \cite{wang2023genemodel}.

\paragraph{Standardisation des variables catégorielles.}
Les variables textuelles (Sexe, Région) ont fait l'objet d'une harmonisation de casse
et d'espaces (\textit{title case}) afin d'éviter les doublons lors de l'encodage.

% -------------------------------------------------------
\subsection{Ingénierie des caractéristiques (Feature Engineering)}
\label{subsec:feateng}
% -------------------------------------------------------

L'objectif du feature engineering est de maximiser le \textbf{Recall} (sensibilité)
du modèle, priorité absolue en diagnostic médical pour minimiser les faux négatifs
\cite{contreras2018ai}. Le tableau~\ref{tab:features} résume les transformations appliquées.

\begin{table}[h!]
\centering
\caption{Récapitulatif des caractéristiques transformées pour le modèle final}
\label{tab:features}
\begin{tabular}{llll}
\hline
\textbf{Catégorie} & \textbf{Variable} & \textbf{Type} & \textbf{Transformation} \\
\hline
Démographique  & Âge au diagnostic             & Numérique  & StandardScaler ($\mu{=}0, \sigma{=}1$) \\
               & Âge (groupe clinique)         & Catégorie  & Discrétisation $<18$, $18$–$30$, $>30$ \\
Génétique      & GRS (67 SNPs)                 & Numérique  & Score de risque pondéré \\
Génomique      & Hub genes (ESCO2, ANP32A-IT1) & Expression & Normalisation RMA / ComBat \\
Immunologique  & Auto-anticorps                & Binaire    & Encodage (Positif=1, Négatif=0) \\
Composée       & Ratio ESCO2/ANP32A-IT1        & Ratio      & Feature synthétique discriminante \\
Catégorielle   & Sexe, Région                  & Nominale   & One-Hot Encoding \\
\hline
\end{tabular}
\end{table}

Le dataset final contient \textbf{1\,000 échantillons} et \textbf{20 features}
après encodage One-Hot. Il a été divisé de façon \textbf{stratifiée} en :

\begin{itemize}
    \item \textbf{Entraînement} : 600 échantillons (60\,\%)
    \item \textbf{Validation} : 200 échantillons (20\,\%)
    \item \textbf{Test} : 200 échantillons (20\,\%)
\end{itemize}

La stratification assure que le taux de prévalence de 25\,\% de cas DT1 est
conservé dans chacun des trois sous-ensembles, évitant tout biais
d'échantillonnage lors de l'évaluation des performances.

% ============================================================
\section{Algorithmes de classification}
\label{sec:algos}
% ============================================================

Six algorithmes ont été entraînés et comparés sur le même jeu de
données : Régression Logistique, Arbre de Décision, Forêt Aléatoire
(\textit{Random Forest}), XGBoost, LightGBM et Support Vector Machine (SVM).
La configuration matérielle utilisée : processeur Intel Core i7 11\textsuperscript{e} génération,
32\,Go de RAM, sans GPU dédié.

% ============================================================
\section{Optimisation des hyperparamètres}
\label{sec:hyperparams}
% ============================================================

L'optimisation a combiné \textit{Grid Search} et \textit{Random Search}
avec validation croisée stratifiée à $K=5$ folds. Les hyperparamètres clés
explorés incluent : \texttt{n\_estimators}, \texttt{max\_depth} et
\texttt{learning\_rate}.

% ============================================================
\section{Métriques d'évaluation}
\label{sec:metriques}
% ============================================================

La métrique prioritaire est le \textbf{Recall} (sensibilité), car minimiser les
faux négatifs est cliniquement critique — un cas de DT1 non détecté peut conduire
à des complications sévères. Les métriques complémentaires utilisées sont :
la Précision, le F1-score et l'AUC-ROC.

% ============================================================
\section{Interprétabilité des modèles}
\label{sec:interpretabilite}
% ============================================================

Les valeurs \textbf{SHAP} (\textit{SHapley Additive exPlanations}) ont été calculées
pour quantifier l'importance globale des features et expliquer individuellement chaque
prédiction aux cliniciens. La méthode \textbf{LIME} (\textit{Local Interpretable
Model-agnostic Explanations}) a été utilisée pour valider les explications locales
par une approche linéaire indépendante.


% Chapitre 4 : Résultats
\chapter{Résultats}
\label{chap:resultats}

\section{Statistiques descriptives}

Distribution des données : 25\% DT1+, 75\% DT1-. Âge moyen DT1+ : 25 ans vs. 35 ans DT1-.
Biomarqueurs : ESCO2 (moyenne DT1+ : 5.24 vs. 2.27), ANP32A-IT1 (4.18 vs. 1.99).

[INSÉRER TABLEAUX ET FIGURES]

\section{Performances des modèles}

\subsection{Comparaison baseline}
[TABLEAU : Accuracy, Recall, Precision, F1, AUC pour les 6 modèles]

\subsection{Modèle final optimisé}
[VOS RÉSULTATS : meilleur modèle, hyperparamètres optimaux, performances test set]

\section{Analyse SHAP}

\subsection{Importance des features}
Top 3 biomarqueurs : [À COMPLÉTER avec vos résultats SHAP]

\subsection{Exemples de prédictions}
Waterfall plots pour patients DT1+ et DT1-.

\section{Application Streamlit}

Démonstration fonctionnelle : upload CSV, prédictions, visualisations, export PDF.


% Chapitre 5 : Discussion
\chapter{Discussion}
\label{chap:discussion}

\section{Validation des biomarqueurs}

Cohérence avec littérature : ESCO2 et ANP32A-IT1 confirment rôles immunitaire et cellulaire \cite{Zhang2022}.
Pertinence pour populations africaines.

\section{Comparaison avec l'état de l'art}

Performances comparables/supérieures aux modèles occidentaux malgré dataset synthétique.
Avantage du focus populations locales.

\section{Limitations de l'étude}

\subsection{Données synthétiques}
Validation sur données réelles camerounaises nécessaire.

\subsection{Taille du dataset}
1000 patients : suffisant pour proof-of-concept, mais échelle clinique requiert plus.

\subsection{Biomarqueurs limités}
Seulement 3 biomarqueurs testés. Autres candidats possibles.

\section{Implications pratiques}

Outil d'aide à la décision pour cliniciens camerounais.
Accessibilité (pas de GPU requis), interface intuitive.

\section{Perspectives}

Validation externe, intégration autres biomarqueurs, études longitudinales.


% Chapitre 6 : Conclusion
\chapter{Conclusion}
\label{chap:conclusion}

\section{Synthèse des contributions}

Ce mémoire a développé un système de détection précoce du DT1 adapté aux populations camerounaises.
Contributions : (1) dataset synthétique réaliste, (2) comparaison 6 algorithmes ML, (3) validation biomarqueurs via SHAP, (4) application démonstrative.

\section{Réponse aux objectifs}

Objectifs atteints : [RÉSUMER vos résultats principaux].
Hypothèses validées/partiellement validées : [BILAN].

\section{Perspectives de recherche}

Court terme : validation données réelles, amélioration recall.
Moyen terme : intégration clinique, formation personnel soignant.
Long terme : déploiement national, extension autres pays africains.

\section{Impact potentiel}

Amélioration diagnostic précoce, réduction mortalité infantile DT1, contribution santé publique camerounaise.

\vspace{1cm}

\begin{center}
\textit{``La science n'a pas de patrie, mais l'homme de science doit penser d'abord aux besoins de son pays."}\\
--- Louis Pasteur
\end{center}


% Bibliographie
\printbibliography[heading=bibintoc, title={Bibliographie}]

% Annexes (optionnelles)
\appendix
\chapter{Code source principal}
\label{ann:code}

Le code source complet du projet est disponible dans l'arborescence du projet, avec les fichiers principaux suivants :

\begin{itemize}
    \item \texttt{2\_DONNEES/generate\_synthetic\_data.py} : Génération des données synthétiques
    \item \texttt{2\_DONNEES/scripts\_preprocessing/01\_nettoyage.py} : Nettoyage des données
    \item \texttt{4\_SCRIPTS\_PRODUCTION/pipeline\_complet.py} : Pipeline ML complet
    \item \texttt{5\_APPLICATION\_DEMO/app\_streamlit.py} : Application de démonstration
\end{itemize}

Les notebooks pédagogiques Jupyter sont disponibles dans le dossier \texttt{3\_NOTEBOOKS\_PEDAGOGIQUES/}.

\chapter{Résultats supplémentaires}
\label{ann:resultats}

[À compléter avec tableaux et figures supplémentaires non inclus dans le corps du mémoire]

\end{document}
